%\chapter*{Введение}
\ssr{Введение}

Цель данной лабораторной работы:
\begin{itemize}
	\item познакомиться с применением структуры непересекающихся множеств и алгоритма Краскала для построения минимального остовного дерева;
	\item изучить эвристические методы поиска пути, реализовать Greedy Best-First Search и A*;
	\item сравнить BFS, алгоритм Дейкстры, Greedy и A* по стоимости найденного маршрута и количеству обработанных состояний.
\end{itemize}

Задачи:
\begin{enumerate}
	\item научиться различать задачу минимального остовного дерева и задачу кратчайшего пути;
	\item научиться использовать DSU с оптимизациями, задавать пространство состояний на клеточном поле;
	\item научиться рассчитывать эвристические оценки;
	\item научиться организовывать OPEN и CLOSED;
	\item научиться восстанавливать маршрут и экспериментально оценивать влияние эвристики на эффективность поиска.
\end{enumerate}

\ssr{Выполнение}

\ssection{Исходные данные}

Исходные граф для задания~1 и клеточное поле для заданий~2-6 представлены в листинге~\ref{og}.

\lstinputlisting[
	label={og},
	style=python, 
	caption={Исходные данные},
	linerange={1-33}
]{main.py}

\ssection{Задание 1. DSU и минимальное остовное дерево}

%\lstinputlisting[
%	label={dsu},
%	style=python, 
%	caption={},
%	linerange={1}
%]{main.py}

% чем задача минимального остовного дерева отличается от поиска кратчайшего пути между двумя вершинами

\ssection{Задание 2. Пространство состояний и базовый BFS}

%\lstinputlisting[
%	label={bfs},
%	style=python, 
%	caption={},
%	linerange={1}
%]{main.py}

\ssection{Задание 3. Эвристические функции}

%\lstinputlisting[
%	label={eurist},
%	style=python, 
%	caption={Функции manhattan(a, b) и euclidean(a, b)},
%	linerange={1}
%]{main.py}

% Сравните полученные значения и объясните, почему для перемещений только по четырём направлениям манхэттенская эвристика обычно информативнее евклидовой

% почему обе эти эвристики на данном поле не завышают минимальную стоимость пути при единичной стоимости каждого перехода

\ssection{Задание 4. Greedy Best-First Search}

%\lstinputlisting[
%	label={gbfs},
%	style=python, 
%	caption={Greedy Best-First Search},
%	linerange={1}
%]{main.py}

% Сравните маршрут Greedy с маршрутом BFS. Объясните, почему ориентация только на h(n) не гарантирует кратчайший путь.

\ssection{Задание 5. Алгоритм A*}

%\lstinputlisting[
%	label={a_asterisk},
%	style=python, 
%	caption={Алгоритм A*},
%	linerange={1}
%]{main.py}

% Сравните две эвристики и сделайте вывод, какая из них дала более направленный поиск на данном поле.

\ssection{Задание 6. Итоговое сравнение BFS, Дейкстры, Greedy и A*}

Итог сравнения алгоритмов представлен в таблице~\ref{compare}.

\begin{table}[!htbp]
	\caption{Сравнение алгоритмов поиска пути}
	\label{compare}
	\small
	\begin{tabularx}{\textwidth}{|X|X|X|X|X|}
		\hline
		\textbf{Алгоритм} & \textbf{Стоимость маршрута} & \textbf{Длина маршрута} & \textbf{Обработано состояний} & \textbf{Гарантия оптимальности} \\\hline
		BFS &&&& при одинаковых весах\\\hline
		Дейкстра &&&& да\\\hline
		Greedy &&&& нет\\\hline
		A* Manhattan &&&& при допустимой эвристике\\\hline
		A* Euclidean &&&& при допустимой эвристике\\\hline	
	\end{tabularx}
\end{table}

% сформулируйте вывод о компромиссе между гарантией оптимальности и количеством обработанных состояний.

\ssr{Вывод}

В ходе лабораторной работы был:
\begin{itemize}
	\item получены знания о применении структуры непересекающихся множеств и алгоритма Краскала для построения минимального остовного дерева;
	\item изучены эвристические методы поиска пути, реализовать Greedy Best-First Search и A*;
	\item проведено сравнение BFS, алгоритм Дейкстры, Greedy и A* по стоимости найденного маршрута и количеству обработанных состояний.
\end{itemize}

\ssr{Контрольные вопросы}
\begin{enumerate}
	\item \textit{Что называется остовным деревом графа?}\newline
	Остовным деревом связного неориентированного графа называется подграф, содержащий все вершины исходного графа, остающийся связным и не содержащий циклов.
	\item \textit{Чем минимальное остовное дерево отличается от кратчайшего пути?}\newline
	Кратчайший путь минимизирует стоимость маршрута между выбранными вершинами, а MST минимизирует суммарную стоимость соединения всех вершин.
	\item \textit{Для чего используется структура DSU?}\newline
	Структура непересекающихся множеств DSU (Disjoint Set Union, Union-Find) поддерживает разбиение элементов на непересекающиеся компоненты. 
	\item \textit{Что выполняют операции find и union?}\newline
	Операция find(x) возвращает представителя компоненты, union(a, b) -- объединяет две компоненты.
	\item \textit{Как работают сжатие путей и объединение по рангу или размеру?}\newline
	
	\item \textit{В чём состоит жадный принцип алгоритма Краскала?}\newline
	Алгоритм Краскала сортирует все рёбра по весу и последовательно добавляет самое легкое ребро, если оно соединяет две разные компоненты.
	\item \textit{Чем принцип работы Прима отличается от Краскала?}\newline
	Алгоритм Прима строит остов от выбранной начальной вершины, каждый раз добавляя самое лёгкое ребро, которое соединяет уже включенную часть дерева с ещё не включённой вершиной.
	\item \textit{Что называется пространством состояний в задаче поиска пути?}\newline
	Задачу поиска маршрута можно представить как поиск в пространстве состояний. 
	Состоянием является положение агента, а ребра графа соответствуют допустимым переходам. 
	\item \textit{Что такое эвристическая функция?}\newline
	Эвристическая функция h(n) -- такая быстро вычисляемая функция. которая отражает направление к цели и  оценивает оставшуюся стоимость пути от текущего состояния n до цели. 
	Она не обязана совпадать с реальной стоимостью пути.
	\item \textit{Чем манхэттенское расстояние отличается от евклидова?}\newline
	Ьанхэттенская оценка обычно точнее отражает минимально необходимое число шагов, тогда как евклидова оценка учитывает прямое геометрическое расстояние.
	\item \textit{Для каких задач может использоваться расстояние Хэмминга?}\newline
	Расстояние Хэмминга подсчитывает количество позиций, в которых два состояния отличаются. ???
	\item \textit{Как Greedy Best-First Search выбирает следующую вершину?}\newline
	Он выбирает из множества OPEN состояние с минимальной эвристической оценкой h(n).
	\item \textit{Почему Greedy не гарантирует оптимальный маршрут?}\newline
	Алгоритм тремится быстрее приблизиться к цели, но не учитывает стоимость уже пройденного пути, поэтому часто рассматривает мало состояний.
	\item \textit{Что означают g(n), h(n) и f(n) в A*?}\newline
	g(n) -- фактическая стоимость пути от старта до n; 
	h(n) -- оценка оставшейся стоимости; 
	f(n) -- суммарный приоритет состояния.
	\item \textit{Для чего в A* используются OPEN и CLOSED?}\newline
	OPEN хранит обнаруженные, но еще не обработанные состояния, обычно в приоритетной очереди. 
	CLOSED содержит состояния, обработка которых завершена.
	\item \textit{Что означает допустимость эвристики?}\newline
	Эвристика называется допустимой, если она никогда не завышает истинную минимальную стоимость пути до цели.
	\item \textit{Что означает согласованность эвристики?}\newline
	Согласованная эвристика дополнительно к дпокстимости требует, чтобы оценка в текущей вершине не превышала стоимость перехода к соседу плюс оценку соседа.
	\item \textit{Во что превращается A* при h(n) = 0?}\newline
	В алгоритм Дейкстры.
	\item \textit{В каком случае BFS гарантирует оптимальный путь?}\newline
	BFS оптимален по числу ребер только при одинаковой стоимости переходов.
	\item \textit{Почему при взвешенных клетках минимальный по числу шагов маршрут может не быть минимальным по стоимости?}\newline
	
	\item \textit{Почему корректная эвристика может уменьшить число обработанных состояний по сравнению с Дейкстрой?}\newline
	
\end{enumerate}